% !TeX program = lualatex
% ================================================================
% Urdu Mathematical Dictionary
%
% Generated by the server using the following settings:
%
%   language            Urdu (urd)
%   text direction      right to left
%   font                Noto Nastaliq Urdu
%   fallback font       Noto Serif
%   built from          latex/dictionary/mathematicalDictionary.tex
%   tier 3 vocabulary   green   RGB 0 128 0
%   English text        black   RGB 0 0 0
%   Urdu text           purple  RGB 112 48 160
%   layout macros       version 2
%
% Please don't edit any information in this comment block - the server
% compares it against the current settings and rebuilds this file when
% they differ, so changes here would be overwritten. However, feel free
% to make updates to the LaTeX below if you want to edit it.
% ================================================================
%
% ---- What is safe to edit in this file ----------------------------------
%
% Safe:     the Urdu word and the Urdu definition - the third
%           and fourth arguments of each entry. That is what this file is for:
%           if a translation reads wrongly to somebody who speaks the language,
%           correct it here and every sheet that uses the word follows.
%
% Not safe: the English word and the English definition - the first and second
%           arguments. The first is how an entry is matched to the shared
%           dictionary; the second is the wording the translation was made
%           from, and is how the server knows the translation is still current.
%           Changing an English definition here does not change the shared
%           dictionary - it only makes this entry look up to date when it is
%           not. Edit the English in the shared dictionary instead, and this
%           file will be brought back into step on the next run.
%
% Entries are added and removed by the server, following the shared dictionary.
% An entry the server cannot read is reported rather than guessed at, so if a
% run complains about this file, look for a missing brace.
% -------------------------------------------------------------------------

\documentclass[11pt]{article}
\usepackage[a4paper,margin=18mm]{geometry}
% Characters this language's font has not got are borrowed from another,
% rather than being silently dropped - see docs/translated-sheet-latex.md
\directlua{%
  if luaotfload and luaotfload.add_fallback then
    luaotfload.add_fallback("eallatinfallback",
      { "Noto Serif:mode=harf;" })
  else
    tex.error("this luaotfload is too old to register a fallback font")
  end
}
\usepackage[english,bidi=basic]{babel}
\babelprovide[import]{urdu}
\babelfont[urdu]{rm}[Renderer=Harfbuzz, BoldFont={Noto Nastaliq Urdu}, ItalicFont={Noto Nastaliq Urdu}, BoldItalicFont={Noto Nastaliq Urdu}, RawFeature={fallback=eallatinfallback}]{Noto Nastaliq Urdu}
\babelfont[urdu]{sf}[Renderer=Harfbuzz, BoldFont={Noto Nastaliq Urdu}, ItalicFont={Noto Nastaliq Urdu}, BoldItalicFont={Noto Nastaliq Urdu}, RawFeature={fallback=eallatinfallback}]{Noto Nastaliq Urdu}
\newcommand{\ealtext}[1]{\foreignlanguage{urdu}{#1}}
\newcommand{\ealtextblock}[1]{{\raggedleft\foreignlanguage{urdu}{#1}\par}}
% ================================================================
% EAL layout helpers
% ================================================================
\usepackage{xcolor}

\definecolor{ealkeycolour}{RGB}{0,128,0}
\definecolor{ealenglish}{RGB}{0,0,0}
\definecolor{ealtwocolour}{RGB}{112,48,160}

\newsavebox{\ealboxen}
\newsavebox{\ealboxtr}
\newlength{\ealglosswd}
\newlength{\ealglossraise}
\setlength{\ealglossraise}{1.15em}

% a tier 3 word, in English and in translation alike
\newcommand{\ealkey}[1]{\textcolor{ealkeycolour}{#1}}
\newcommand{\ealkeytr}[1]{\textcolor{ealkeycolour}{\ealtext{#1}}}

% One translated word above one English word. Built from kernel
% primitives, never a tabular, which would break wherever the sheet
% loads array, booktabs or tabularx. The box is as wide as the wider of
% the two, so a long translation pushes its neighbours aside instead of
% overlapping them, and it claims its full height so a table row or a
% TikZ node makes room for it.
\newcommand{\ealgloss}[2]{%
  \sbox{\ealboxen}{\ealkey{#1}}%
  \sbox{\ealboxtr}{\small\textcolor{ealtwocolour}{\ealtext{#2}}}%
  \setlength{\ealglosswd}{\wd\ealboxen}%
  \ifdim\wd\ealboxtr>\ealglosswd\setlength{\ealglosswd}{\wd\ealboxtr}\fi
  \leavevmode
  \makebox[\ealglosswd][c]{%
    \makebox[0pt][c]{\usebox{\ealboxen}}%
    \makebox[0pt][c]{\raisebox{\ealglossraise}{\usebox{\ealboxtr}}}%
  }%
}

% opens up line spacing around a block of glosses, so the raised
% translations sit clear of the line above
\newenvironment{ealglossed}{\par\linespread{2.1}\selectfont\ignorespaces}{\par}

% a whole translated sentence above its English counterpart. block level,
% so it does not belong inside a tabular cell
\newcommand{\ealpara}[2]{%
  \par\addvspace{0.5\baselineskip}%
  {\color{ealtwocolour}\ealtextblock{#1}}%
  \penalty200
  {\color{ealenglish}#2\par}%
  \addvspace{0.5\baselineskip}%
}

% ================================================================
% Dictionary layout
% ================================================================
\usepackage{multicol}

\newlength{\ealdictgap}
\setlength{\ealdictgap}{1.2mm}

% english word / english meaning / translated word / translated meaning
\newcommand{\dictentrytr}[4]{%
  \par\addvspace{\ealdictgap}%
  \noindent
  {\bfseries\ealkey{#1}}\quad{\small #2}\par
  \nopagebreak
  {\ealtextblock{{\bfseries\color{ealkeycolour}#3}~\textemdash{} {\small #4}}}%
  \par
}

\pagestyle{empty}
\setlength{\parindent}{0pt}

\begin{document}
{\large\bfseries Urdu Mathematical Dictionary}\par\medskip

\dictentrytr{coefficient}{the number in front of a letter, showing how many of it there are}{عددی سر}{کسی حرف سے پہلے لکھا ہوا عدد، جو ظاہر کرتا ہے کہ اس حرف کی کتنی تعداد موجود ہے}
\dictentrytr{decimal}{a number written with a point, where the digits after the point are parts of a whole}{اعشاریہ}{ایک عدد جو ایک نقطے کے ساتھ لکھا جاتا ہے، جہاں نقطے کے بعد کے ہندسے ایک پورے کے حصے ہوتے ہیں}
\dictentrytr{denominator}{the number below the line in a fraction, showing how many equal parts the whole is split into}{مخرج}{کسر کی لکیر کے نیچے والا عدد، جو ظاہر کرتا ہے کہ پورا کتنے برابر حصوں میں تقسیم ہوا ہے}
\dictentrytr{difference}{the answer when one number is subtracted from another}{فرق}{وہ جواب جو ایک عدد کو دوسرے عدد سے منہا کرنے پر حاصل ہو}
\dictentrytr{digit}{a single figure from 0 to 9 that a number is written with}{ہندسہ}{0 سے 9 تک کا ایک اکیلا عددی نشان، جس سے کوئی عدد لکھا جاتا ہے}
\dictentrytr{equivalent}{having the same value, even though it is written differently}{مساوی}{ایک ہی قدر رکھنے والا، چاہے وہ مختلف طریقے سے لکھا گیا ہو}
\dictentrytr{estimate}{a sensible answer worked out quickly from rounded numbers}{تخمینہ}{ایک معقول جواب جو مقرب اعداد سے جلدی سے حاصل کیا جائے}
\dictentrytr{expand}{to multiply out brackets so the expression has no brackets left}{توسیع کرنا}{قوسین کو ضرب دے کر کھولنا تاکہ اظہار میں کوئی قوس باقی نہ رہے}
\dictentrytr{expression}{a group of numbers and letters joined by operations, with no equals sign}{اظہار}{اعداد اور حروف کا ایک گروہ جو عملیات سے جڑا ہوا ہو، جس میں علامتِ مساوات نہ ہو}
\dictentrytr{factor}{a whole number that divides exactly into another number}{عامل}{ایک مکمل عدد جو کسی دوسرے عدد کو بغیر باقی کے تقسیم کرے}
\dictentrytr{factorise}{to write an expression as a product, usually by putting brackets back in}{تجزیه کرنا}{کسی اظہار کو حاصل ضرب کے طور پر لکھنا، عموماً قوسین واپس لگا کر}
\dictentrytr{formula}{a rule written with letters, showing how to work one quantity out from others}{فارمولا}{ایک قاعدہ جو حروف سے لکھا جاتا ہے، جو بتاتا ہے کہ ایک مقدار کو دوسری مقداروں سے کیسے حاصل کیا جائے}
\dictentrytr{fraction}{a number written as one whole split into equal parts, with a number of those parts taken}{کسر}{ایک عدد جو اس طرح لکھا جاتا ہے کہ ایک پورا برابر حصوں میں تقسیم ہو اور ان حصوں میں سے کچھ حصے لیے گئے ہوں}
\dictentrytr{gradient}{how steep a line is, found by how far it rises for each step across}{میلان}{وہ مقدار جو ظاہر کرتی ہے کہ ایک لکیر کتنی ڈھلوان ہے، اور جو اس بات سے معلوم ہوتی ہے کہ ہر اکائی افقی فاصلے پر وہ کتنی اوپر اٹھتی ہے}
\dictentrytr{highest common factor}{the largest whole number that divides exactly into two or more numbers}{سب سے بڑا مشترک عامل}{سب سے بڑا مکمل عدد جو دو یا زیادہ اعداد کو بغیر باقی کے تقسیم کرے}
\dictentrytr{identity}{a statement that is true for every value of the letters in it}{سنجاک}{ایک بیان جو اس میں موجود حروف کی ہر قدر کے لیے درست ہو}
\dictentrytr{index}{the small raised number showing what power something is raised to}{اس}{وہ چھوٹا اوپر والا عدد جو ظاہر کرتا ہے کہ کسی عدد کو کس قوت تک بڑھایا گیا ہے}
\dictentrytr{inequality}{a statement that one quantity is larger or smaller than another}{نامساوات}{ایک بیان کہ ایک مقدار دوسری مقدار سے بڑی یا چھوٹی ہے}
\dictentrytr{integer}{a whole number, which may be positive, negative or zero}{صحيح عدد}{ایک مکمل عدد، جو مثبت، منفی یا صفر ہو سکتا ہے}
\dictentrytr{lowest common multiple}{the smallest number that two or more numbers all divide into exactly}{سب سے چھوٹا مشترک مضاعف}{سب سے چھوٹا عدد جسے دو یا زیادہ اعداد بغیر باقی کے تقسیم کریں}
\dictentrytr{multiple}{a number you get by multiplying another number by a whole number}{مضاعف}{ایک عدد جو کسی دوسرے عدد کو ایک مکمل عدد سے ضرب دینے سے حاصل ہو}
\dictentrytr{multiply}{to add a number to itself a given number of times}{ضرب دینا}{کسی عدد کو ایک دی گئی تعداد میں خود سے جمع کرنا}
\dictentrytr{numerator}{the number above the line in a fraction, showing how many equal parts are taken}{صورت}{کسر کی لکیر کے اوپر والا عدد، جو ظاہر کرتا ہے کہ کتنے برابر حصے لیے گئے ہیں}
\dictentrytr{percentage}{a number of parts out of a hundred}{فیصد}{سو میں سے کچھ حصوں کی تعداد}
\dictentrytr{place value}{the value a digit has because of where it sits in a number}{مقامی قدر}{وہ قدر جو ایک ہندسے کو عدد میں اس کی جگہ کی وجہ سے حاصل ہوتی ہے}
\dictentrytr{power}{how many times a number is multiplied by itself}{قوت}{وہ تعداد جو ظاہر کرتی ہے کہ کسی عدد کو خود سے کتنی بار ضرب دیا جاتا ہے}
\dictentrytr{prime}{a whole number greater than 1 that divides only by itself and 1}{مفرد عدد}{1 سے بڑا مکمل عدد جو صرف اپنے آپ اور 1 سے تقسیم ہو سکے}
\dictentrytr{product}{the answer when two or more numbers are multiplied}{حاصل ضرب}{وہ جواب جو دو یا زیادہ اعداد کو ضرب دینے سے حاصل ہو}
\dictentrytr{proportion}{a relationship where two amounts change by the same factor}{تناسب}{ایک تعلق جس میں دو مقداریں ایک ہی عامل سے تبدیل ہوتی ہیں}
\dictentrytr{quotient}{the answer when one number is divided by another}{خارجِ قسمت}{وہ جواب جو ایک عدد کو دوسرے عدد سے تقسیم کرنے سے حاصل ہو}
\dictentrytr{ratio}{a way of comparing two or more amounts, showing how much of one there is for each of the other}{نسبت}{دو یا زیادہ مقداروں کا موازنہ کرنے کا ایک طریقہ، جو ظاہر کرتا ہے کہ ایک مقدار دوسری مقداروں میں سے ہر ایک کے مقابلے میں کتنی ہے}
\dictentrytr{reciprocal}{the number you multiply by to get 1, found by turning a fraction upside down}{مقلوب}{وہ عدد جس سے ضرب دینے پر 1 حاصل ہو، جو کسر کو الٹا کرنے سے ملتا ہے}
\dictentrytr{recur}{to repeat the same digits forever after the decimal point}{تکرار پذیر ہونا}{اعشاریہ نقطے کے بعد وہی ہندسے ہمیشہ کے لیے دہرائے جانا}
\dictentrytr{root}{a number that gives another number when multiplied by itself a given number of times}{جذر}{ایک عدد جو خود سے ایک مقررہ تعداد میں ضرب دینے پر دوسرا عدد دیتا ہے}
\dictentrytr{round}{to replace a number with a nearby simpler one}{مقرب کرنا}{کسی عدد کو اس کے قریب کے آسان عدد سے بدل دینا}
\dictentrytr{scale factor}{the number every length is multiplied by when a shape is enlarged}{پیمانہ عامل}{وہ عدد جس سے کسی شکل کو بڑا کرتے وقت ہر لمبائی کو ضرب دیا جاتا ہے}
\dictentrytr{simplify}{to write something in its smallest or plainest form without changing its value}{سادہ کرنا}{کسی چیز کو اس کی سب سے چھوٹی یا سادہ ترین شکل میں لکھنا، اس کی قدر کو تبدیل کیے بغیر}
\dictentrytr{sum}{the answer when numbers are added together}{حاصل جمع}{وہ جواب جو اعداد کو ایک ساتھ جمع کرنے سے حاصل ہو}
\dictentrytr{surd}{a root that cannot be written exactly as a fraction, so it is left in root form}{جذر اصم}{ایک جذر جسے کسر کی صورت میں بالکل نہ لکھا جا سکے، اس لیے اسے جذر کی شکل میں چھوڑ دیا جاتا ہے}
\dictentrytr{unitary}{a method that first finds the value of one part or one item}{اکائی طریقہ}{ایک طریقہ جو پہلے ایک حصے یا ایک چیز کی قدر معلوم کرتا ہے}

\end{document}
